\section{Universal equations}\label{sec:equations}

The completion does not change the universal equational theory of the free
completion in the language obtained by naming the elements of the base.

\begin{theorem}[Preservation and reflection of equations]
\label{thm:equational-conservativity}
For every algebraic equation $s=t$ in variables, the basic operation symbols,
and constants naming elements of $A$,
\[
 \Free D\models s=t
 \quad\Longleftrightarrow\quad
 \FCcomp{\Free D}\models s=t.
\]
\end{theorem}

\begin{proof}
The reverse implication follows by restricting an equation valid in the
completion to the embedded copy of $\Free D$.

For the forward implication, suppose $\Free D\models s=t$.  Each basic
operation preserves every congruence $\fceq{m}$ and therefore extends
continuously to the Cauchy completion; the same is true of every term
operation.  Given an assignment of the variables in the completion and a fixed
$m$, choose representatives from $\Free D$ sufficiently close to the assigned
values at level $m$.  Evaluating $s$ and $t$ on these representatives gives the
same element because the equation holds in $\Free D$.  Continuity of term
evaluation then shows that the two values of the completed terms are
$\fceq{m}$-equivalent.  Since this holds for every $m$ and the uniformity is
separated, the two completed values are equal.
\end{proof}

No corresponding claim is made here for arbitrary conditional equations; the
argument above uses only universal equational validity and continuity of term
operations.
